% 张三慧V1第11章热力学第二定律候选独立解析。
% 原题物理PDF398--400；配套书物理PDF178--184。
% 热/功符号及第11.13题原书与辅导书的等熵、等焓疑点均需专家核查。
\section*{第11章\quad 热力学第二定律：独立解析}

\paragraph{习题 11.1}
图11.14的折线 \(a\to b\to c\) 由
\(a=(1\,\mathrm L,6\cdot10^5\,\mathrm{Pa})\)、
\(b=(2\,\mathrm L,8\cdot10^5\,\mathrm{Pa})\)、
\(c=(3\,\mathrm L,4\cdot10^5\,\mathrm{Pa})\) 构成。
两段为直线，做功
\[
 W_{ab}=\frac{p_a+p_b}{2}(V_b-V_a)=700\,\mathrm J,
 \qquad W_{bc}=\frac{p_b+p_c}{2}(V_c-V_b)=600\,\mathrm J.
\]
故 \(\boxed{W_{ac}=1300\,\mathrm J}\)。
一摩尔单原子氦气
\(U=\tfrac32pV\)，因
\(p_aV_a=600\,\mathrm J\)、
\(p_cV_c=1200\,\mathrm J\)，
\(\Delta U=900\,\mathrm J\)。
第一定律得到
\(\boxed{Q_{ac}=2200\,\mathrm J}\)。
熵为状态函数，按状态方程可沿任意可逆辅助路径积分：
\[
 \boxed{\Delta S
 =\frac32R\ln\frac{T_c}{T_a}+R\ln\frac{V_c}{V_a}
 =R\left(\frac32\ln2+\ln3\right)
 \approx17.8\,\mathrm{J/K}}.
\]
不可用实际折线的 \(Q_{ac}/T_{\rm avg}\)
来代替熵积分。

\paragraph{习题 11.2}
令冰起始 \(T_i=233.15\,\mathrm K\)，
熔点 \(T_m=273.15\,\mathrm K\)，
沸点 \(T_b=373.15\,\mathrm K\)。
四阶段吸热是冰升温、冰熔化、水升温、
水汽化：
\[
 \boxed{Q=m[c_i(T_m-T_i)+\lambda
 +c_w(T_b-T_m)+L]}.
\]
使用 \(m=30\,\mathrm g\)，得到
\[
 Q=30[2.1(40)+334+4.2(100)+2260]
 =\boxed{9.294\times10^4\,\mathrm J}.
\]
把不可逆的实际加热以微元平衡辅助过程计算熵，
\[
 \boxed{\Delta S
 =m\left[c_i\ln\frac{T_m}{T_i}
 +\frac{\lambda}{T_m}
 +c_w\ln\frac{T_b}{T_m}
 +\frac L{T_b}\right]}.
\]
这里的潜热比热单位为 \(\mathrm{J/g}\)，
结果约 \(3.16\times10^2\,\mathrm{J/K}\)。
水蒸气被视为100℃刚汽化完成。

\paragraph{习题 11.3}
环境近似为温度
\(T_{\rm amb}=273\,\mathrm K\) 的恒温库。
人体每日传给环境热量
\(Q=8\times10^5\,\mathrm J\)，
故环境熵增加
\[
 \boxed{\Delta S_{\rm amb}=\frac Q{T_{\rm amb}}
 \approx2.93\times10^3\,\mathrm{J/K\ per\ day}}.
\]
该值不是人体本身的熵变，也未计进食或呼吸。

\paragraph{习题 11.4}
由室内高温侧
\(T_h=293.15\,\mathrm K\)
向室外低温侧
\(T_c=253.15\,\mathrm K\) 传热，
热功率
\(\dot Q=2\times10^8/3600\,\mathrm W\)。
室内熵变化率 \(-\dot Q/T_h\)，
室外为 \(+\dot Q/T_c\)，
整个系统生成熵率
\[
 \boxed{\dot S_{\rm gen}=\dot Q
 \left(\frac1{T_c}-\frac1{T_h}\right)
 \approx30.0\,\mathrm{J/(K\,s)}>0}.
\]

\paragraph{习题 11.5}
匀速车摩擦中机械功率不可逆转化为热
\(\dot Q=20\,\mathrm{kJ/s}\)，
环境温度 \(T=285.15\,\mathrm K\)：
\[
 \boxed{\dot S_{\rm gen}=\frac{\dot Q}{T}
 =\frac{20000}{285.15}
 \approx70.1\,\mathrm{J/(K\,s)}}.
\]
只看稳态外部摩擦热源的熵产生，
不把机械有序能先算成热熵。

\paragraph{习题 11.6}
每秒水质量
\(\dot m=\rho\dot V
 =(1000)(1500)=1.5\times10^6\,\mathrm{kg/s}\)。
落差 \(h=74\,\mathrm m\)，
重力势能流率
\[
 \dot E_g=\dot m gh
 \approx(1.5\times10^6)(9.80)(74)
 =1.0878\times10^9\,\mathrm W.
\]
温度12℃近似285.15K，
故
\[
 \boxed{\dot S_{\rm gen}=\frac{\dot E_g}{T}
 \approx3.81\times10^6\,\mathrm{J/(K\,s)}}.
\]
瀑布宽度83m未进入计算，
因为总流量已经给出。

\paragraph{习题 11.7}
设水比热容
\(C=mc=4200\,\mathrm{J/K}\)，
\(T_0=273.15\,\mathrm K\)，
\(T_{50}=323.15\,\mathrm K\)，
\(T_{100}=373.15\,\mathrm K\)。
水升温的熵增始终是
\[
 \Delta S_w=C\ln(T_{100}/T_0).
\]
(1) 只有100℃热库时，热库放热
\(Q=C(T_{100}-T_0)\)，
系统总熵变
\[
 \boxed{\Delta S_{\rm tot,1}
 =C\ln\frac{T_{100}}{T_0}
 -C\frac{T_{100}-T_0}{T_{100}}>0}.
\]
(2) 两级热库先50℃再100℃，
熵变
\[
 \boxed{\Delta S_{\rm tot,2}
 =C\ln\frac{T_{100}}{T_0}
 -C\frac{T_{50}-T_0}{T_{50}}
 -C\frac{T_{100}-T_{50}}{T_{100}}>0}.
\]
由于加热过程分两级，温差缩小，
\(\Delta S_{\rm tot,2}<\Delta S_{\rm tot,1}\)；
两热库各损失的熵要分别算，
不能将两热源视为同温。

\paragraph{习题 11.8}
初态2.5kg水、0.7kg冰，总质量3.2kg，
目标第一次达到水、冰\emph{等质量}，
即各1.6kg。
因此\emph{有0.9kg水凝固}，
而不是0.7kg冰全部熔化。
取水熔化潜热
\(\lambda=3.34\times10^5\,\mathrm{J/kg}\)，
熔点 \(T_m=273.15\,\mathrm K\)。
(1) 向略低于冰点的房间放热
\(Q=0.9\lambda\)，
\[
 \boxed{\Delta S_{\rm mix,1}
 =-\frac{0.9\lambda}{T_m}
 \approx-1.10\times10^3\,\mathrm{J/K}},
\]
\[
 \boxed{\Delta S_{\rm total,1}
 =0.9\lambda\left(\frac1{T_{\rm room}}
 -\frac1{T_m}\right)>0
 \quad(T_{\rm room}<T_m)}.
\]
源书“比0℃稍有不同”未指定精确值，
故总熵只可给含 \(T_{\rm room}\) 的表达式。
(2) 放到100℃的恒温热库，
熔解0.9kg冰恢复原始相态，
\[
 \boxed{\Delta S_{\rm mix,2}
 =+\frac{0.9\lambda}{T_m}},
\]
\[
 \boxed{\Delta S_{\rm total,2}
 =0.9\lambda\left(\frac1{T_m}
 -\frac1{373.15}\right)>0}.
\]
\emph{由热传递导致的相变是不可逆的}，
但熵变计算可以借助熔点处的可逆辅助路径。

\paragraph{习题 11.9}
初态 \(p_1V_1=(3.039\times10^5)(4)
 =1.2156\times10^6\,\mathrm J\)，
\(T_1=300\,\mathrm K\)，
设 \(C_V=nR/(\gamma-1)\) 为整份气体定容热容。
第一段从4到16立方米等温：
\[
 \boxed{W_{12}=Q_{12}
 =p_1V_1\ln4\approx1.685\times10^6\,\mathrm J},
\quad
 \boxed{\Delta S_{12}=\frac{p_1V_1}{T_1}\ln4}.
\]
第二段等容降温直到其末态3能够沿可逆绝热
返回初态，故
\[
 T_3=T_1(4/16)^{\gamma-1}
 =300\,4^{-0.4}\approx172.3\,\mathrm K,
\]
\[
 p_3=p_1(4/16)^\gamma,\qquad
 W_{23}=0,\quad Q_{23}=\Delta U_{23}
 =C_V(T_3-T_1)<0,
\]
\[
 \boxed{\Delta S_{23}
 =C_V\ln(T_3/T_1)=-\Delta S_{12}}.
\]
第三段可逆绝热压缩从16到4立方米：
\[
 Q_{31}=\Delta S_{31}=0,\quad
 \boxed{W_{31}=C_V(T_3-T_1)<0}.
\]
注意 \(\Delta U_{31}=C_V(T_1-T_3)>0\)。
一周结束 \(\boxed{\Delta S_{\rm cycle}=0,
 \Delta U_{\rm cycle}=0}\)，
\[
 \boxed{W_{\rm cycle}=Q_{\rm cycle}
 =p_1V_1\ln4+C_V(T_3-T_1)}.
\]
路径是等温向右、等容向下、绝热向左上返回。

\paragraph{习题 11.10}
绝热容器内两份等质量同种液体，
常数比热 \(c_p\)。
由能量守恒
\(mc_pT_1+mc_pT_2=2mc_pT_f\)，
\[
 \boxed{T_f=(T_1+T_2)/2}.
\]
熵变用可逆辅助升降温计算：
\[
 \boxed{\Delta S=mc_p\ln\frac{T_f}{T_1}
 +mc_p\ln\frac{T_f}{T_2}
 =mc_p\ln\frac{(T_1+T_2)^2}{4T_1T_2}}.
\]
由 \((T_1-T_2)^2\ge0\) 得
\((T_1+T_2)^2\ge4T_1T_2\)，
故 \(\Delta S\ge0\)；
当 \(T_1\ne T_2\) 时严格大于零。

\paragraph{习题 11.11}
两份各 \(\nu\) mol同种理想气体，
初温相同为 \(T\)，
压强 \(p_1,p_2\)。
各自容积 \(V_i=\nu RT/p_i\)，
刚性绝热连通后，理想气体内能只由T决定，
因此终温仍为T，
总容积 \(V_1+V_2\)，
终压
\(p_f=2\nu RT/(V_1+V_2)
 =2p_1p_2/(p_1+p_2)\)。
可用各份初态到终态的等温可逆辅助过程算熵：
\[
 \Delta S=\nu R\ln(p_1/p_f)
 +\nu R\ln(p_2/p_f),
\]
化简
\[
 \boxed{\Delta S=\nu R
 \ln\frac{(p_1+p_2)^2}{4p_1p_2}\ge0}.
\]
由算术几何平均不等式得等号条件 \(p_1=p_2\)；
两份气体相同不能另加虚假的“混合熵”。

\paragraph{习题 11.12}
金属与液体同处恒压绝热环境，
忽略相变，比热视常数。
令 \(C_1=mc_p,C_2=m'c'_p\)，
能量守恒给
\[
 \boxed{T_f=\frac{C_1T_i+C_2T'_i}{C_1+C_2}}.
\]
系统熵变
\[
 \boxed{\Delta S=C_1\ln\frac{T_f}{T_i}
 +C_2\ln\frac{T_f}{T'_i}\ge0}.
\]
证明最终态是\emph{固定总能量约束下的熵极大态}：
对两部分末态温度 \(T,T'\) 与约束
\(C_1T+C_2T'=E\)，
总熵一阶变分为
\(dS=C_1dT/T+C_2dT'/T'\)。
用 \(C_1dT+C_2dT'=0\) 得
\(dS=C_1(1/T-1/T')dT\)，
极值必有 \(T=T'\)。
沿约束变量的二阶导
\(-C_1/T^2-C_1^2/(C_2T'^2)<0\)，
故为最大值，不只是驻点。

\paragraph{习题 11.13}
本题\emph{必须以主教材对热力过程的定义为准}。
主教材V1 PDF物理页399给定
(1)(2)是节流膨胀，(3)是可逆绝热膨胀。
理想稳流节流满足\emph{等焓}
\(H_m=\mathrm{const}\)，
准静态可逆绝热满足\emph{等熵}
\(S_m=\mathrm{const}\)。
应在图11.16氮气温熵图中
分别沿相应\emph{等焓线}或
\emph{竖直等熵线}到指定的终压等压线，
然后读末温。
配套辅导书S PDF物理页184
列出近似温度
\[
 \boxed{T_1\simeq95\,\mathrm K,\quad
 T_2\simeq65\,\mathrm K,\quad
 T_3\simeq20\,\mathrm K}.
\]
\textbf{重大源书差异：}
配套书对前两问的过程称为“等熵”，
第三问称“等焓”，
其文字与主教材给的
“节流(1)(2)、可逆绝热(3)”\emph{反向}。
故上列三个温度只能先登记为
\emph{配套答案图读值}，
不能声称已经通过主教材指定路径的独立核验。
第三问的20K还处于图示氮气液汽共存区，
不能把它简单代回理想气体绝热公式；
本题待专门重绘图11.16、按正确路径逐点重读。
